\chapter{Relato de Experiência}
\label{cap:relato}

\section{Contextualização do Projeto}

Explique inicialmente  por que o software envolvido no projeto é relevante: que problema ele resolve, para que público-alvo e com que nível de criticidade para as partes interessadas no software.

Explique, em seguida, o cenário de cada projeto em que atuou: qual problema existia, por que o software foi criado ou modificado no projeto, ou seja, qual a finalidade do projeto, e quem são as partes interessadas no projeto.

Sintetize, ainda, a \textbf{metodologia de desenvolvimento ou manutenção de software} usada pela equipe de Engenharia de Software do projeto (Scrum, Kanban etc.). 

\subsection{Ambiente Organizacional}

Descreva o ambiente organizacional em que o projeto foi realizado, em termos de cultura organizacional, abordagem de processos (definidos ou informais, rigorosos ou ágeis), local e infraestrutura de trabalho, pessoas envolvidas e seus respectivos papeis, forma de trabalho (presencial, remoto ou híbrido), dimensão da equipe, entre outras características relevantes para compreender o ambiente em que o projeto se desenvolveu.


Explique sucintamente como o projeto foi organizado em termos de \textbf{metodologia de desenvolvimento ou manutenção de software} usada pela equipe de Engenharia de Software do projeto (Scrum, Kanban, XP, RUP, etc.), mostrando os macro-processos.

Discuta sucintamente as principais ferramentas, tecnologias e \textit{frameworks} (Git, Figma, Docker, etc.) aplicados no projeto.

Descreva, ainda, os papéis e responsabilidades da equipe, principalmente aqueles que tinham interface ou interação com o papel que você exerceu.

\subsection{Forma de Participação no Projeto}

Apresente o contexto geral da sua participação no(s) projeto(s) de software. Retome o tipo de projeto (desenvolvimento, evolução, adaptação, correção, migração, etc.) e o seu(s) papel(is) na equipe envolvida. 

Explique como você passou a integrar o(s) projeto(s) e quais eram as expectativas iniciais em relação às atividades a serem desempenhadas.

Explique, ainda, como você se familiarizou com o projeto antes de começar a contribuir efetivamente de acordo com o seu papel no projeto.

\section{Atividades Realizadas e Resultados Gerados}

Detalhe as atividades específicas que você desempenhou ao longo do projeto, tais como engenharia de requisitos, modelagem, projeto de software, implementação, testes, documentação de processos, manutenção corretiva ou evolutiva, entre outras. 

A descrição deve evidenciar como cada atividade foi atribuída a você e como foi conduzida por você: quais ferramentas ou técnicas foram utilizadas, resultados gerados e como estes contribuíram para o avanço do projeto.

Descreva o que foi feito, como foi feito e por quê.

\section{Integração com a Equipe e Processo de Trabalho}

Apresente a forma como você interagiu com os demais membros da equipe. Discuta aspectos como comunicação, colaboração, reuniões, uso de ferramentas de gestão e práticas adotadas pelo grupo. 

Descreva também como o processo de trabalho da equipe influenciou as suas atividades, considerando metodologias, rotinas e dinâmicas internas.

\section{Desafios Enfrentados}

Discuta os principais desafios encontrados durante a sua experiência na Fábrica de Software. Os desafios podem ser das mais diversas naturezas, tais como técnicos, organizacionais, comunicacionais, relacionados ao aprendizado de novas ferramentas e práticas ou ao desenvolvimento de habilidades pessoais e comportamentais. 

A análise deve destacar como esses desafios impactaram o andamento das suas atividades e quais estratégias foram utilizadas para superá-los.

\section{Soluções Adotadas e Resultados Obtidos}

Discuta as soluções implementadas para lidar com os desafios identificados e os resultados alcançados com as atividades realizadas. 
Os resultados podem incluir entregas concluídas, melhorias no processo, correções implementadas, artefatos produzidos, habilidades adquiridas ou aprimoradas, ou qualquer outro impacto positivo decorrente da sua participação no(s) projeto(s).

Analise o impacto do trabalho para as partes interessadas: o que melhorou, o que funcionou, o que não funcionou.
Analise criticamente a evolução do produto e a avaliação das partes interessadas, tanto usuários quanto outros \textit{stakeholders}, em comparação com os objetivos iniciais do projeto.

\section{Avaliação de Qualidade de Software}

Apresente os conceitos relacionados à qualidade de software das atividades que você realizou, incluindo critérios de qualidade, indicadores e métricas usadas na avaliação, estratégias de verificação e validação e práticas de garantia de qualidade aplicadas ao trabalho realizado no TCC. 

Fundamente a qualidade dos artefatos que você gerou durante o TCC com base em evidências produzidas por atividades como testes, revisão de código, \textit{walkthrough} de documentação ou qualquer outra ação voltada à avaliação da qualidade do produto ou do processo de software.

\section{Aprendizados e Desenvolvimento Profissional}

Discuta os aprendizados adquiridos ao longo da experiência, dos pontos de vista técnico, profissional e pessoal. 

Devem ser abordados conhecimentos consolidados, habilidades (\textit{hard} e \textit{soft}) desenvolvidas, boas práticas assimiladas e reflexões sobre o trabalho em equipe e sobre o processo de engenharia de software. 

Analise, com destaque, se, e como, a experiência realizada na Fábrica de Software contribuiu para a sua formação integral como profissional e como cidadão.
